\section[Stage 031]{Stage 031: Reachability and Stable-Side Monotonicity}
\label{stage:031}

\stagefield{Part / anchor}{Part II; primary anchor \resultanchor{MTDC-T5}.}
\claimstatus{\StatusExactClosure{} inside the selected lower-mode branch.}

\stagefield{Purpose}{Stage~031 proves that the selected prefactor is strictly increasing as the rank-one loading approaches the softening threshold.  This turns selected normalization into a one-dimensional reachability theorem.}

\stagefield{Inputs}{The stage uses \(\lambda_-(\alpha_0)\), \(s_-(\alpha_0)\), \(P_{0,-}=\beta_0s_-/\lambda_-\), and the stable interval \(0\leq\alpha_0<\alpha_{\rm crit}\).}

\stagefield{Verification}{SymPy audit: \StageFile{scripts/moving_throat_pde_stage031_selected_branch_reachability_sympy_audit.py}; transcript: \StageFile{scripts/output/moving_throat_pde_stage031_selected_branch_reachability_sympy_audit.txt}.  Mathematica audit: \StageFile{mathematica/moving_throat_pde_stage031_selected_branch_reachability_mathematica_audit.wl}; transcript: \StageFile{mathematica/output/moving_throat_pde_stage031_selected_branch_reachability_mathematica_audit.txt}.}

\subsection*{Derivation}

\paragraph{1. Derivative of the selected overlap.}
The exact derivative is
\begin{equation}
\label{eq:app-stage031-ds-dalpha}
\boxed{
\frac{ds_-}{d\alpha_0}=
\frac{2(\Delta K_{\rm ax})^2\Pi_\kappa}{R_\lambda^3}>0}.
\end{equation}
Using \(d\lambda_-/d\alpha_0=-s_-\), one obtains
\begin{equation}
\label{eq:app-stage031-dP-dalpha}
\boxed{
\frac{dP_{0,-}}{d\alpha_0}
=\beta_0\frac{(ds_-/d\alpha_0)\lambda_-+s_-^2}{\lambda_-^2}>0
}
\end{equation}
whenever \(\lambda_->0\) and \(\beta_0>0\).

\paragraph{2. Endpoints.}
At zero loading,
\begin{equation}
\label{eq:app-stage031-P-zero}
P_{0,-}(0)=\frac{\beta_0\kappa_0^2}{A}.
\end{equation}
The refined threshold, now with the isotropic \(U\)-softening absorbed into \(A,B\), is
\begin{equation}
\label{eq:app-stage031-alpha-crit}
\boxed{
\alpha_{\rm crit}=\frac{AB}{B\kappa_0^2+A\kappa_1^2}}.
\end{equation}
As \(\alpha_0\to\alpha_{\rm crit}^-\), \(\lambda_-\to0^+\) while \(s_-\) remains positive, so
\begin{equation}
P_{0,-}\to+\infty.
\end{equation}
Thus the stable branch spans
\begin{equation}
\label{eq:app-stage031-range}
P_{0,-}(0)\leq P_{0,-}<+\infty.
\end{equation}

\paragraph{3. Unique crossing.}
Because \(P_{0,-}\) is continuous and strictly increasing on the stable interval, every target satisfying
\begin{equation}
P_{\rm target}>P_{0,-}(0)
\end{equation}
has a unique \(\alpha_*\in(0,\alpha_{\rm crit})\) such that \(P_{0,-}(\alpha_*)=P_{\rm target}\).  For the quadrupole target,
\begin{equation}
P_{\rm target}=\frac{N_Q^{\rm target}}{\widehat m_-^{\,2}}.
\end{equation}

\stagefield{Output}{Stage~031 outputs the derivative identities \eqref{eq:app-stage031-ds-dalpha}--\eqref{eq:app-stage031-dP-dalpha}, the refined threshold \eqref{eq:app-stage031-alpha-crit}, and the unique stable crossing theorem.}

\stagefield{Checks}{\begin{verificationchecklist}
\item The positivity of \(ds_-/d\alpha_0\) follows from \((\Delta K_{\rm ax})^2\Pi_\kappa>0\) and \(R_\lambda>0\).
\item The derivative of \(P_{0,-}\) is positive because both terms in the numerator are nonnegative and \(s_-^2>0\).
\item The divergence at softening follows from \(\lambda_-\to0^+\) with \(s_->0\).
\end{verificationchecklist}}

\stagefield{Downstream use}{Stage~032 modifies the target by fixing \(\widehat m_-\).  Stage~035 upgrades the reachability statement into the explicit monotone shape function \(F(\xi,\delta)\).}

\clearpage

